\documentclass[11pt,a4paper]{article}

% --- Essential Packages ---
\usepackage[utf8]{inputenc}
\usepackage[T1]{fontenc}
\usepackage{amsmath}   % Required for \text{} inside math mode
\usepackage{amssymb}   % For advanced mathematical symbols
\usepackage{geometry}  % To cleanly manage page margins
\geometry{margin=1in}

\begin{document}
	
	Good evening! That is a beautifully intuitive way to think about it. Heating a metal sheet gives us a perfect physical canvas for differential geometry. 
	
	To prove that the gradient of your temperature function $f$ transforms as a \textbf{covariant vector} (or a dual vector/one-form), we need to look at how its components change when you change your coordinate system. 
	
	The punchline of the proof is this: \textbf{a covariant vector transforms using the partial derivatives of the \textit{old} coordinates with respect to the \textit{new} ones.} This is a direct consequence of the multivariable \textbf{chain rule}.
	
	Here is the step-by-step proof.
	
	\medskip
	\hrule
	\medskip
	
	\section*{The Proof}
	
	\subsection*{1. Setting up two coordinate systems}
	Imagine you describe the points on your metal sheet using a standard Cartesian coordinate system, $x^\mu = (x^1, x^2)$. 
	
	Now, imagine a colleague decides to look at the same sheet using a different coordinate system, say polar coordinates or just a shifted/rotated grid, which we will call $\tilde{x}^\nu = (\tilde{x}^1, \tilde{x}^2)$. 
	
	The temperature at any physical point on the sheet is an intrinsic physical reality---it doesn't care what numbers you write on your grid. Therefore, $f$ is a \textbf{scalar field}, meaning its value at a given point is invariant under a change of coordinates:
	\[
	\tilde{f}(\tilde{x}) = f(x)
	\]
	
	\subsection*{2. Taking the Gradient}
	In your coordinate system, the components of the gradient are given by the partial derivatives:
	\[
	\nabla_\mu f = \frac{\partial f}{\partial x^\mu}
	\]
	
	In your colleague's coordinate system, the components are:
	\[
	\nabla_\nu \tilde{f} = \frac{\partial \tilde{f}}{\partial \tilde{x}^\nu}
	\]
	
	\subsection*{3. Applying the Chain Rule}
	To see how your colleague's gradient components relate to yours, we apply the multi-variable chain rule to express the derivative with respect to $\tilde{x}^\nu$:
	\[
	\frac{\partial \tilde{f}}{\partial \tilde{x}^\nu} = \frac{\partial f}{\partial x^\mu} \frac{\partial x^\mu}{\partial \tilde{x}^\nu}
	\]
	\textit{(Note: We are using Einstein summation notation here, where the repeated index $\mu$ implies a sum over $\mu = 1, 2$.)}
	
	\subsection*{4. Analyzing the Transformation Law}
	Let's rewrite that result by substituting our gradient components:
	\[
	\nabla_\nu \tilde{f} = \frac{\partial x^\mu}{\partial \tilde{x}^\nu} \nabla_\mu f
	\]
	
	Look closely at the Jacobian matrix element $\frac{\partial x^\mu}{\partial \tilde{x}^\nu}$. Because the ``new'' coordinates $\tilde{x}$ are in the \textit{denominator}, this is the exact mathematical definition of a \textbf{covariant transformation law}. 
	
	By definition, a geometric object whose components transform like:
	\[
	V_\nu' = \frac{\partial x^\mu}{\partial x'^\nu} V_\mu
	\]
	is a covariant vector. Therefore, the gradient $\nabla f$ is a covariant vector field.
	
	\medskip
	\hrule
	\medskip
	
	\section*{Physical Intuition: Why ``Co''-variant?}
	
	It helps to contrast this with regular positional vectors (like a displacement vector $dx^\mu$), which are \textbf{contravariant}. 
	
	If you decide to shrink the scale of your grid lines (e.g., measuring in millimeters instead of centimeters), the \textit{numerical values} of your coordinates change. 
	
	\begin{itemize}
		\item \textbf{Contravariant (e.g., Distance):} The distance numbers get \textbf{larger} ($1\text{ cm}$ becomes $10\text{ mm}$). The components vary \textit{contrary} to the change of scale.
		\item \textbf{Covariant (e.g., Gradient):} The temperature gradient measures \textit{change per unit distance} ($^\circ\text{C}/\text{cm}$). If you switch to millimeters, the temperature change per millimeter gets \textbf{smaller} ($1^\circ\text{C}/\text{cm}$ becomes $0.1^\circ\text{C}/\text{mm}$). The components vary \textit{together} with the scale change.
	\end{itemize}
	
	Because the gradient components live naturally as coefficients to displacements (forming the invariant scalar $df = \frac{\partial f}{\partial x^\mu} dx^\mu$), the gradient inherently belongs to the dual space, acting as a machine that eats contravariant vectors to output a scalar temperature change.
	
\end{document}